% SPDX-License-Identifier: MIT
% Copyright (c) 2026 Evgeny Baulin
% Optimization Methods, Seminar 04: one-page cheat sheet (Russian).
% Generated by checks/04. Accelerated Gradient Methods.ipynb from web/04/js/content-cheatsheet.js and content/students.json. Do not edit by hand.
% Compile with: latexmk -xelatex -interaction=nonstopmode -halt-on-error Cheatsheet_ru.tex
\RequirePackage{fix-cm}
\documentclass[10pt,a4paper]{article}

\usepackage{amsmath,amssymb}
\usepackage{fontspec}
\setmainfont{cmunrm.otf}[BoldFont=cmunbx.otf, ItalicFont=cmunti.otf, BoldItalicFont=cmunbi.otf]
\setsansfont{cmunss.otf}[BoldFont=cmunsx.otf, ItalicFont=cmunsi.otf, BoldItalicFont=cmunso.otf]
\setmonofont{cmuntt.otf}[BoldFont=cmuntb.otf, ItalicFont=cmunit.otf, BoldItalicFont=cmuntx.otf]
\usepackage{polyglossia}
\setdefaultlanguage{russian}
\setotherlanguage{english}
\newfontfamily\cyrillicfont{cmunrm.otf}[Ligatures=TeX, BoldFont=cmunbx.otf, ItalicFont=cmunti.otf, BoldItalicFont=cmunbi.otf]
\newfontfamily\cyrillicfontsf{cmunss.otf}[Ligatures=TeX, BoldFont=cmunsx.otf, ItalicFont=cmunsi.otf, BoldItalicFont=cmunso.otf]
\newfontfamily\cyrillicfonttt{cmuntt.otf}[BoldFont=cmuntb.otf, ItalicFont=cmunit.otf, BoldItalicFont=cmuntx.otf]
\PolyglossiaSetup{russian}{hyphenmins={2,3}}

\usepackage[a4paper,left=8mm,right=8mm,top=7mm,bottom=11mm,footskip=5.5mm]{geometry}
\usepackage{xcolor}
\usepackage{multicol}
\usepackage{enumitem}
\usepackage{ragged2e}
\usepackage{fancyhdr}
\usepackage[colorlinks=true,linkcolor=hseblue,urlcolor=hseblue]{hyperref}

\definecolor{hseblue}{RGB}{16,45,105}
\definecolor{hseaccent}{RGB}{198,40,40}
\definecolor{hsegray}{RGB}{108,114,126}
\definecolor{hselight}{RGB}{240,243,248}

\hypersetup{pdftitle={Шпаргалка к семинару 4: Ускоренные градиентные методы}, pdfauthor={Евгений Баулин}}

% Page style: the author and the license in the footer.
\pagestyle{fancy}
\fancyhf{}
\renewcommand{\headrulewidth}{0pt}
\fancyfoot[L]{\scriptsize\color{hsegray}\textcopyright\ 2026 Евгений Баулин, \href{https://t.me/tarakan_tuc}{t.me/tarakan\_tuc}}
\fancyfoot[R]{\scriptsize\color{hsegray}Лицензия MIT}

% Compact layout.
\setlength{\parindent}{0pt}
\setlength{\parskip}{0pt}
\setlength{\columnsep}{4.5mm}
\setlength{\columnseprule}{0.3pt}
\renewcommand{\columnseprulecolor}{\color{hselight}}
\setlength{\emergencystretch}{2em}
\setlength{\RaggedRightRightskip}{0pt plus 4em}
\binoppenalty=2000
\relpenalty=500
\setlength{\abovedisplayskip}{2pt}
\setlength{\belowdisplayskip}{2pt}
\setlist[itemize]{leftmargin=7pt, labelsep=3pt, itemsep=1.2pt, topsep=1pt, parsep=0pt, partopsep=0pt,
  label={\color{hseaccent}\scriptsize\textbullet}}

\makeatletter
\newcommand{\cssec}[1]{\par\addvspace{3.5pt}\noindent
  \colorbox{hselight}{\parbox{\dimexpr\linewidth-2\fboxsep}{\RaggedRight\hyphenpenalty=10000\color{hseblue}\bfseries\strut #1}}%
  \par\nobreak\vspace{1pt}\@afterheading}
\makeatother
\clubpenalty=10000
\widowpenalty=10000

\begin{document}
\thispagestyle{fancy}

\noindent
\begin{minipage}[b]{0.6\textwidth}
  \RaggedRight\hyphenpenalty=10000
  {\color{hseaccent}\bfseries\small Шпаргалка к семинару 4}\\[1pt]
  {\color{hseblue}\bfseries\Large Ускоренные\\ градиентные методы}
\end{minipage}\hfill
\begin{minipage}[b]{0.39\textwidth}
  \raggedleft\footnotesize
  {\bfseries Методы оптимизации}\\
  НИУ ВШЭ, факультет компьютерных наук, 2026--2027\\
  Лектор курса: Андрей Игнатов\\
  Автор: Евгений Баулин, Telegram \href{https://t.me/tarakan_tuc}{@tarakan\_tuc}
\end{minipage}

\vspace{2pt}
{\color{hseblue}\rule{\textwidth}{0.8pt}}
\vspace{-4pt}

\fontsize{9}{10.6}\selectfont
\begin{multicols}{3}
\RaggedRight


% ---- quadratic
\cssec{Квадратичная функция и GD}
\begin{itemize}
  \item $f=\tfrac{1}{2}x^\top Qx-b^\top x$, $x^*=Q^{-1}b$, $\kappa=L/\mu$.
  \item $0<\alpha<2/L$; лучший постоянный квадратичный шаг $2/(L+\mu)$.
\end{itemize}

% ---- hb
\cssec{Тяжёлый шар}
\begin{itemize}
  \item $x^+=x-\alpha g(x)+\beta(x-x^-)$; $x_{-1}=x_0$.
  \item $p_\lambda(r)=r^2-(1+\beta-\alpha\lambda)r+\beta$.
  \item При $0\le\beta<1$: $0<\alpha<2(1+\beta)/L$.
  \item $\alpha_*={4\over(\sqrt L+\sqrt\mu)^2}$, $\beta_*=({\sqrt L-\sqrt\mu\over\sqrt L+\sqrt\mu})^2$.
  \item $\rho_*=({\sqrt\kappa-1})/({\sqrt\kappa+1})$. Результат для квадратичной функции; возможен начальный рост.
\end{itemize}

% ---- nag
\cssec{Опережение}
\begin{itemize}
  \item $y=x+\beta(x-x^-)$, $x^+=y-\alpha g(y)$.
  \item Градиент в экстраполированной точке. Рестарт отбрасывает инерцию и повторяет шаг.
\end{itemize}

% ---- cd
\cssec{Координатный спуск}
\begin{itemize}
  \item $L_i=Q_{ii}$, $x_i^+=x_i-g_i/L_i$. Уменьшение $g_i^2/(2L_i)$.
  \item Жадный выбор: максимум $|g_i|/\sqrt{L_i}$. Циклический и равномерный случайный выбор — другие правила.
  \item $d$ обновлений = один проход. Учитывайте стоимость жадного просмотра.
\end{itemize}

% ---- cg
\cssec{CG для положительно определённой системы}
\begin{itemize}
  \item Вначале $r=b-Ax$, $p=r$. Переиспользуйте $q=Ap$.
  \item $\alpha={r^\top r\over p^\top q}$; $x^+=x+\alpha p$, $r^+=r-\alpha q$.
  \item $\beta={(r^+)^\top r^+\over r^\top r}$; $p^+=r^++\beta p$.
  \item $r_i^\top r_j=0$, $p_i^\top Ap_j=0$ в точной арифметике, $i\ne j$.
  \item Проверяйте истинную $b-Ax$; относительная невязка использует $\|b\|$, если она ненулевая.
\end{itemize}

% ---- sgd
\cssec{SGD и дисперсия}
\begin{itemize}
  \item $F=N^{-1}\sum_i\ell_i$, $g_B=|B|^{-1}\sum_{i\in B}\nabla\ell_i$.
  \item Независимые выборки: $\operatorname{Var}(g_B)=\sigma^2/B$.
  \item Скалярная квадратичная функция: $V_\infty={\alpha\sigma^2\over B\mu(2-\alpha\mu)}$, $0<\alpha\mu<2$.
  \item В начале эпохи создавайте новую перестановку; включайте последний короткий батч и правильно нормируйте его.
\end{itemize}

% ---- schedule
\cssec{Расписания и работа}
\begin{itemize}
  \item $\alpha_0$, $\alpha_0/\sqrt{1+k}$, $\alpha_0/(1+k)$; $k$ считает обновления.
  \item Сравнивайте $5$ seed, центральную кривую и интервал между квартилями. Считайте градиенты объектов и вычисления полной потери.
\end{itemize}

% ---- scale
\cssec{Масштаб признаков}
\begin{itemize}
  \item $Z=(X-m)S^{-1}$, $H_Z=S^{-1}HS^{-1}$.
  \item Исходные коэффициенты $w=S^{-1}u$, константа $b=c-m^\top w$.
  \item Используйте обучающие статистики. Масштабирование не устраняет корреляцию.
\end{itemize}

\end{multicols}
\end{document}
